\subsection*{B.5 Entrevista N.°~5 --- Shakira (After Trash)}
\addcontentsline{toc}{subsection}{B.5 Entrevista N.° 5 --- Shakira (After Trash)}

\noindent\textbf{Entrevistado:} Shakira (nombre real)\\
\noindent\textbf{Perfil:} Dueña de After Trash, emprendimiento de indumentaria estampada (remeras, buzos, tote bags con diseños de bandas, películas y cultura pop). Tienda virtual activa desde marzo--abril de 2025; local físico abierto hace aproximadamente dos meses (mayo de 2026). Modelo de producción bajo pedido (\textit{make-to-order}).\\
\noindent\textbf{Entrevistador:} Tomás Gonçalves (moderador)\\
\noindent\textbf{Fecha:} Domingo 27 de julio de 2026\\
\noindent\textbf{Modalidad:} Videollamada (semiestructurada) --- Duración: \textasciitilde35 minutos\\
\noindent\textbf{Plataforma del negocio:} Tienda Negocio (\textasciitilde\$20.000/mes, sin comisión por venta). Anteriormente en Tienda Nube.

\begin{framed}
\small\noindent\textit{Nota metodológica: Esta transcripción fue limpiada a partir de una grabación de \textasciitilde35 minutos vía Google Meet. Se eliminó charla trivial, se corrigió la sintaxis para mejorar la fluidez de lectura y se suprimieron algunas repeticiones. No se agregó información que la entrevistada no haya dicho explícitamente.}
\end{framed}

\subsubsection*{Bloque 0 · Introducción y rapport}

\noindent\textbf{Tomás Gonçalves:} Nosotros estamos haciendo nuestro proyecto final: una herramienta que ayuda a dueños de tienda online a tomar mejores decisiones usando sus propios datos. Nosotros le decimos recomendaciones accionables. La idea es que agarra directamente tus datos y en base a eso te dice, por ejemplo: el martes que viene podés poner esta promoción y tenés asegurada esta cantidad de ventas.

\noindent\textbf{Tomás Gonçalves:} Para empezar: ¿qué vendés y cómo se llama tu emprendimiento?

\noindent\textbf{Shakira:} Soy dueña de After Trash. Vendo indumentaria estampada: remeras, buzos, tote bags. Los diseños son de bandas de rock, películas, series, cosas de cultura pop en general.

\noindent\textbf{Tomás Gonçalves:} ¿Desde cuándo estás?

\noindent\textbf{Shakira:} Empecé entre marzo y abril de 2025. Este año abrí el local, hace como dos meses.

\subsubsection*{Bloque 1 · Contexto del negocio y herramientas actuales}

\noindent\textit{--- Plataforma de venta y migración ---}

\noindent\textbf{Tomás Gonçalves:} ¿Estás en Tienda Nube?

\noindent\textbf{Shakira:} No, estoy en Tienda Negocio. Antes estaba en Tienda Nube pero la migré. Tienda Negocio me cobra alrededor de veinte mil pesos por mes y no me cobra comisión por venta. Tienda Nube me cobraba comisión además del abono mensual.

\noindent\textbf{Tomás Gonçalves:} ¿Y la diferencia en funcionalidades?

\noindent\textbf{Shakira:} Es muy similar. La diferencia es que Tienda Nube tiene más integraciones y más apps disponibles. Tienda Negocio es más básica pero me alcanza.

\noindent\textit{--- Modelo de producción bajo pedido ---}

\noindent\textbf{Tomás Gonçalves:} ¿Cómo manejás el stock? ¿Producís antes o producís cuando te piden?

\noindent\textbf{Shakira:} Produzco cuando me piden, o sea que es bajo pedido. En la tienda online dejo los productos activos sin límite de stock porque técnicamente puedo hacer cualquier cantidad. Eso simplifica mucho el tema del inventario: no hay riesgo de que me quede sin stock ni de que me sobre mercadería.

\noindent\textbf{Tomás Gonçalves:} ¿Cuánto tarda la producción una vez que recibís el pedido?

\noindent\textbf{Shakira:} Depende de la carga de trabajo que tenga. Normalmente entre tres y cinco días hábiles. En épocas de mucha demanda, como el Hot Sale, puede ser un poco más.

\noindent\textit{--- Registro de ventas y uso de ChatGPT ---}

\noindent\textbf{Tomás Gonçalves:} ¿Cómo llevás el registro de ventas?

\noindent\textbf{Shakira:} La plataforma me da un historial de órdenes. Aparte de eso, uso mucho ChatGPT. Es como si fuese mi socio: le paso datos y me ayuda a analizarlos, a redactar mensajes para clientes, a pensar estrategias. Lo uso constantemente.

\noindent\textbf{Tomás Gonçalves:} ¿Le pasás los datos de ventas directamente?

\noindent\textbf{Shakira:} Sí. Le copio las órdenes o le paso un resumen y le pregunto cosas: cuánto vendí este mes, cuál es el producto que más sale, en qué días vendí más. Me ayuda a procesar la información que tengo dispersa.

\subsubsection*{Bloque 2 · Gestión de inventario y stock}

\noindent\textit{--- Stock infinito online vs. local físico ---}

\noindent\textbf{Tomás Gonçalves:} El local físico, ¿tiene stock separado de la tienda online?

\noindent\textbf{Shakira:} Sí. En el local tengo unidades físicas. En la tienda online el stock es infinito porque produzco bajo pedido. Eso hace que no pueda sincronizar bien los dos canales: si alguien compra en el local una remera de determinado diseño, eso no descuenta nada en la tienda online porque online ese producto siempre figura disponible.

\noindent\textbf{Tomás Gonçalves:} ¿Eso te genera algún problema?

\noindent\textbf{Shakira:} Por ahora no mucho. El local es nuevo y la mayoría de las ventas siguen siendo online. Pero capaz que a futuro sí se convierte en un problema si alguien compra algo online que tengo en exhibición en el local y termino teniendo que producir otro.

\subsubsection*{Bloque 3 · Análisis de ventas y decisiones comerciales}

\noindent\textit{--- Incidente de precios: 35 millones vendidos al costo ---}

\noindent\textbf{Tomás Gonçalves:} ¿Alguna vez tuviste un problema con los precios?

\noindent\textbf{Shakira:} Sí, uno grande. Vendí treinta y cinco millones de pesos esencialmente al costo, sin ganancia. Me pasó porque no actualicé los precios a tiempo. El dólar subió, los insumos aumentaron y yo seguía vendiendo a los precios viejos. Me di cuenta tarde.

\noindent\textbf{Tomás Gonçalves:} ¿Cómo lo detectaste?

\noindent\textbf{Shakira:} Cuando fui a hacer el balance del mes y no me cerraban los números. Vendí mucho ese período pero al calcular el margen real vi que prácticamente no había ganado nada. Fue un aprendizaje muy caro.

\noindent\textbf{Tomás Gonçalves:} ¿Con qué frecuencia actualizás los precios ahora?

\noindent\textbf{Shakira:} Cada vez que hay un cambio importante en los insumos o en el tipo de cambio. No tengo una frecuencia fija, pero estoy más atenta. Si pudiera tener una alerta automática que me diga ``tus precios no se actualizaron en X días y el costo de los insumos subió'', eso sería ideal.

\noindent\textit{--- Estacionalidad y ventas por eventos ---}

\noindent\textbf{Tomás Gonçalves:} ¿Hay épocas que vendés más?

\noindent\textbf{Shakira:} Sí, muy marcadas. El Hot Sale y el Black Friday son los dos picos más grandes. También cuando hay eventos relacionados con los artistas que tengo en los diseños: si Indio Solari da un show, se disparan las ventas de todo lo que tenga que ver con Los Redondos. Es algo muy difícil de predecir de antemano.

\noindent\textbf{Tomás Gonçalves:} ¿Preparás algo antes de esos eventos o te agarran de sorpresa?

\noindent\textbf{Shakira:} Para el Hot Sale y el Black Friday me preparo: pongo descuentos, genero contenido en redes. Para los eventos de artistas es más reactivo; cuando se anuncia algo intento subirme a la ola rápido, publicar en redes y aprovechar el momento.

\noindent\textit{--- Redes sociales y publicidad ---}

\noindent\textbf{Tomás Gonçalves:} ¿Hacés publicidad paga?

\noindent\textbf{Shakira:} Sí, en Instagram y TikTok. La mayoría de los clientes me llegan por redes sociales. También tengo bastante tráfico orgánico porque posteo mucho contenido.

\noindent\textbf{Tomás Gonçalves:} ¿Sabés de dónde viene cada venta?

\noindent\textbf{Shakira:} Más o menos. La plataforma me dice de dónde viene el tráfico pero no siempre es claro. Hay ventas que no tienen fuente asignada. Lo que sí noto es que cuando hago una publicación que pega bien, al día siguiente hay más ventas.

\subsubsection*{Bloque 4 · Puntos de dolor y visibilidad del negocio}

\noindent\textbf{Tomás Gonçalves:} ¿Cuál es el problema más grande que tenés hoy en la gestión del negocio?

\noindent\textbf{Shakira:} El tema de los precios es el que más me preocupa ahora, después de lo que me pasó. Necesito tener una forma de saber cuándo mis precios están desactualizados antes de que me pase de nuevo. El otro problema es no saber bien qué productos tienen más demanda potencial: qué diseños le gustarían a la gente que todavía no hice.

\noindent\textbf{Tomás Gonçalves:} ¿Cómo decidís qué diseños hacer?

\noindent\textbf{Shakira:} Por intuición y por lo que veo en redes. Veo qué está trending, qué artistas están activos, qué le pregunta la gente en comentarios. Pero es muy empírico, no tengo datos.

\subsubsection*{Bloque 5 · Validación de la propuesta de valor}

\noindent\textit{--- Confianza y disposición a pagar ---}

\noindent\textbf{Tomás Gonçalves:} Si tuvieras una herramienta que todos los días te dice cosas como ``tus precios de remeras no se actualizaron en 30 días, el costo del insumo principal subió un 12\%'', ``mañana es una fecha de alto tráfico, activá una promoción'', ``el diseño X tuvo muchas visitas pero pocas ventas, revisá el precio o la descripción'', ¿cuánto pagarías?

\noindent\textbf{Shakira:} Si me ayuda a no perder plata como la vez que te conté, pagaría bastante. Con una prueba gratis para ver que funciona, podría pagar hasta treinta mil pesos por mes. Pero necesito ver que realmente sirve antes de comprometerme.

\noindent\textbf{Tomás Gonçalves:} ¿Un período de prueba te convencería?

\noindent\textbf{Shakira:} Sí, totalmente. Si en el período de prueba me muestra algo útil que yo no estaba viendo, ahí me convence. La herramienta sola no alcanza; necesito ver el valor en la práctica.

\noindent\textit{--- Uso de IA y actitud hacia la tecnología ---}

\noindent\textbf{Tomás Gonçalves:} Ya que usás mucho ChatGPT, ¿cómo ves el uso de IA en el negocio?

\noindent\textbf{Shakira:} Lo veo como algo positivo. ChatGPT me cambió la manera de trabajar. Antes tardaba mucho más en hacer cosas: responder mensajes, analizar datos, redactar publicaciones. Ahora es mucho más rápido. Si hay una herramienta de IA específica para e-commerce que haga algo parecido pero más especializada, bienvenida.

\noindent\textbf{Tomás Gonçalves:} ¿Y si la herramienta te da una recomendación que no te cierra? ¿La seguís igual?

\noindent\textbf{Shakira:} No, la cuestiono. Con ChatGPT también pasa: a veces me da algo que no tiene sentido para mi negocio y lo desestimo. La herramienta tiene que poder explicar por qué recomienda lo que recomienda; si no entiendo el razonamiento, no lo aplico.

\subsubsection*{Bloque 6 · Cierre}

\noindent\textbf{Tomás Gonçalves:} ¿Hay algo que no te pregunté y que te parece importante?

\noindent\textbf{Shakira:} Una cosa que te diría es que la interfaz tiene que ser simple. Yo no soy una persona técnica; sé usar las herramientas que uso porque son intuitivas. Si la herramienta requiere mucha configuración o es confusa, la abandono rápido.

\noindent\textbf{Shakira:} Y lo otro: que tenga en cuenta el contexto argentino. Los precios cambian mucho acá, el tipo de cambio es un factor constante. Una herramienta pensada para mercados estables no me va a servir igual.

\noindent\textbf{Tomás Gonçalves:} Perfecto, eso es muy valioso. Mil gracias, Shakira.

\begin{center}\textit{--- Fin de la entrevista ---}\end{center}
